% Glorious technicolour

\newcommand{\kindc}[1]{{\color{ForestGreen}{#1}}} % c for colour
\newcommand{\typec}[1]{{\color{RoyalBlue}{#1}}}
\newcommand{\termc}[1]{{\color{RedOrange}{#1}}}
\newcommand{\blk}[1]{{\color{black}{#1}}}
\newcommand{\gray}[1]{{\color{gray}{#1}}}
\newcommand{\pic}[1]{{\color{ACMLightBlue}{#1}}}

% abbreviation of the calculus
\newcommand\algst{AlgST\xspace}
\newcommand\fuse{FuSe$^{\{\}}$\xspace}

% Grammars
\newcommand{\grmeq}{\; ::= \;\;}
\newcommand{\grmor}{\;\mid\;}

% KEYWORDS
\newcommand{\syntax}{\texttt}
\newcommand{\Keyword}[1]{\mathsf{#1}}

\newcommand{\link}{\Keyword{lin}}
\newcommand{\unk}{\Keyword{un}}

% KINDS
% \newcommand{\kindstyle}[1]{\kindc{\textsc{\textbf{\upshape #1}}}}
\newcommand{\kindstyle}[1]{\kindc{\textbf{\textsc{#1}}}}
\newcommand{\kindm}{\kindstyle m}
\newcommand{\kindp}{\kindstyle p}
\newcommand{\kinds}{\kindstyle s}
\newcommand{\kindf}{\kindstyle f}
\newcommand{\kindl}{\kindstyle l}
\newcommand{\kindg}{\kindstyle g}
\newcommand{\kindt}{\kindstyle t}
\newcommand{\kind}{\kindc{\kappa}}
% \newcommand{\kind}{\kindc K}
\newcommand{\prekind}{\kindc J}
\newcommand{\kinda}[2]{\kindc{\overline{#1}}\rightarrow\kindc{#2}}
\newcommand{\kindas}[2]{\kindc{#1}\rightarrow\kindc{#2}} % simple, no overline

% MULTIPLICITY
\newcommand\mult{\kindc{m}}
\newcommand\lin{\kindstyle{lin}}
\newcommand\un{\kindstyle{un}}

\newcommand{\kindtu}{\kindt^\un}
\newcommand{\kindtl}{\kindt^\lin}
\newcommand{\kindtm}{\kindt^\mult}
\newcommand{\kindsu}{\kinds^\un}
\newcommand{\kindsl}{\kinds^\lin}
\newcommand{\kindmu}{\kindm^\un}
\newcommand{\kindml}{\kindm^\lin}
\newcommand{\kindfu}{\kindf^\un}
\newcommand{\kindfl}{\kindf^\lin}

\newcommand{\kindgu}{\kindg^\un}
\newcommand{\kindgl}{\kindg^\lin}
\newcommand{\kindpu}{\kindp^\un}
\newcommand{\kindpl}{\kindp^\lin}

% OPERATORS
\newcommand{\operator}[1]{\blk{\operatorname{#1}}}
\newcommand\klub{\sqcup}
\newcommand{\polarizedparam}[2]{\blk{#1}\blk(\typec{#2}\blk)}
% \newcommand{\polarizeminus}[2]{\tosession{\polarizedparam - {#1}}{#2}} % Use \TFIn
% \newcommand{\polarizeplus}[2]{\tosession{\polarizedparam + {#1}}{#2}}  % Use TFOut
\newcommand{\tosessionop}{\mathsection}
\newcommand{\tosession}[2]{\blk{\tosessionop(}\typec{#1}\blk{).}\typec{#2}} % imaginative syntax
%\newcommand{\tosession}[2]{\blk{\tosessionop(}\typec{#1}\blk,\typec{#2}\blk)} % explicit binary
\newcommand\ppos[1]{\typec{+#1}} % was {1^+}
\newcommand\pneg[1]{\typec{-#1}} % was {#1^-}
\newcommand\ppm[1]{\typec{\pm#1}}

% TYPES
\newcommand{\types}[1]{\typec{\syntax{#1}}} % s for syntax
\newcommand\TEndW{\ensuremath{\types{End}_{\typec ?}}} % wait
\newcommand\TEndT{\ensuremath{\types{End}_{\typec !}}} % terminate
\newcommand\TEnd{\types{End}}
\newcommand{\TG}[2]{\typec{{#1}.{#2}}} % DEPRECATED, no guards anymore
\newcommand{\TIn}[2]{\TG{?#1}{#2}}
\newcommand{\TOut}[2]{\TG{!#1}{#2}}
\newcommand{\TGIn}[1]{\typec{?#1}}
\newcommand{\TGOut}[1]{\typec{!#1}}
\newcommand{\TSkip}{\types{Skip}}
\newcommand{\TSeq}[2]{\typec{#1;#2}}
\newcommand{\TChoice}[2]{\typec{#1\oplus#2}}
\newcommand{\TBranch}[2]{\typec{#1\mathbin{\&} #2}}
\newcommand{\TOne}{\types{1}}
\newcommand\TFIn[2]{\tosession{\polarizedparam - {#1}}{#2}}
\newcommand\TFOut[2]{\tosession{\polarizedparam + {#1}}{#2}}
\newcommand\TFMinus[1]{\blk{-(}\typec{#1}\blk)}
\newcommand\protocolk{\types{protocol}}
\newcommand\datak{\types{data}}
\newcommand{\TUp}[1]{\typec{\ensuremath{{\uparrow}{#1}}}}
\newcommand{\TUnit}{\types{Unit}}
\newcommand{\TInt}{\types{Int}}
\newcommand{\TBool}{\types{Bool}}
\newcommand\TFun[3][\mult]{\typec{#2 \to_{#1} #3}}
\newcommand{\TPair}[2]{\typec{#1 \otimes #2}}
% \newcommand{\TPair}[2]{\typec{\langle #1, #2\rangle}}
\newcommand{\TSum}[2]{\typec{#1 + #2}}
\newcommand{\TRecord}[2]{\typec{\{\overline{\termc{#1}\colon{#2}}\}}}
\newcommand{\TVariant}[2]{\typec{[\overline{\termc{#1}\colon{#2}}]}}
\newcommand{\dualk}{\types{Dual}}
% \newcommand{\dualinnerk}{\operator{dual'}}
\newcommand{\TDual}[1]{\typec{\dualk~{#1}}}
% \newcommand{\TDualInner}[1]{\dualinnerk~\typec{#1}}
\newcommand{\TMinus}[1]{\typec{-{#1}}}
\newcommand{\tbind}[2]{\typec{{#1}\blk\colon\kindc{#2}}}
\newcommand{\TForall}[3]{\typec{\forall\tbind{#1}{#2}{.#3}}}
% \newcommand{\TForall}[3]{\typec{\forall(\tbind{#1}{#2}){\,#3}}}
\newcommand{\TBForall}[2]{\typec{\forall{#1}{#2}}}
% \newcommand{\TMForall}[2]{\typec{\forall{#1}{#2}}} # multiple
\newcommand{\TPApp}[2]{\typec{{#1}~{#2}}}
\newcommand{\matchin}[1][i]{\TFIn   {\overline{T_{#1}}\tsubs{\overline U}{\overline\alpha}}{S}}
\newcommand{\matchout}[1][i]{\TFOut {\overline{T_{#1}}\tsubs{\overline U}{\overline\alpha}}{S'}}
\newcommand{\protocols}{\protocolk\ \typec{\proto\
    \overline{\alpha}\colon\overline\kindg}~\blk{=}~\typec{\overline{\constr\
      \overline{\ppm T}}}
}
% \newcommand{\protocolsfullobsolete}{\protocolk\ \typec{\proto\
%     (\overline{\alpha}\colon \overline\kind)}~\blk{=}~\typec{\{{\constr_i : \overline{P_i}}\}_{i\in I}}
% }
\newcommand{\protocolsfull}{\protocolk\ \typec{\proto\ \overline{\alpha}}~\typec{=}~\typec{\{{\termc{\constr_i}\ \overline{T_i}}\}_{i\in I}}
}

\newcommand{\protocolsdataobsolete}{
  \datak\ \typec{D\
    (\overline{\alpha}\colon\overline{\kindc{\kindt}})}~\blk{=}~
  \typec{\{{\constr_i : \overline{T_i}}\}_{i\in I}}
}
\newcommand{\protocolsdata}{
  \datak\ \typec{D\ \overline{\alpha}}~\blk{=}~
  \typec{\{{\constr_i\ \overline{T_i}}\}_{i\in I}}
}

\newcommand\PRep{\types{Rep}}
\newcommand\PSeq{\types{Seq}}
\newcommand\PEps{\types{Eps}}
\NewDocumentCommand\PDispatch{o}{\typec{\types{Dispatch}\IfValueT{#1}{[#1]}}}

% EXPRESSIONS
\newcommand{\exps}[1]{\termc{\syntax{#1}}}

\newcommand{\casek}{\exps{case}}
\newcommand{\ofk}{\exps{of}}
\newcommand{\sendIntk}{\exps{sendInt}}
\newcommand{\receiveIntk}{\exps{receiveInt}}
\newcommand{\sendk}{\exps{send}}
%\newcommand{\selectk}{\exps{select}} % Use C_i alone
\newcommand{\receivek}{\exps{receive}}
\newcommand{\offerk}{\exps{offer}}
\newcommand{\letk}{\exps{let}}
\newcommand{\ink}{\exps{in}}
\newcommand{\newk}{\exps{new}}
\newcommand{\reck}{\exps{rec}}
\newcommand{\forkk}{\exps{fork}}
\newcommand{\matchk}{\exps{match}}
\newcommand{\withk}{\exps{with}}
\newcommand{\terminatek}{\exps{terminate}}
\newcommand{\waitk}{\exps{wait}}
\newcommand{\unitk}{\exps{*}} % {\exps{unit}}
\newcommand{\selectk}{\exps{select}}

% \newcommand{\ESelect}[1]{\termc{\selectk\,{#1}}} % Use C_1 alone
\newcommand{\constr}{\termc C} % constructor
\newcommand{\protodata}{\typec X} % protocoland constructor
% protocol names
\newcommand{\proto}{\typec{\ensuremath{\rho}}}%{\typec P} % protocol
\newcommand{\data}{\typec D} % dtaa

\newcommand{\abse}[3]{\termc{\lambda{#1}\colon{\typec #2}.{#3}}}
\newcommand{\absne}[3]{\termc{\lambda{#1}.{#3}}}
\newcommand{\tabse}[3]{\termc{\Lambda}\,\typec{{#1}\colon{#2}}\termc{. #3}}
\newcommand{\rece}[3]{\termc{\reck\ {#1}\colon{\typec #2}.{#3}}}
% \newcommand{\rece}[5]{\appe{\appe{\tappe{\tappe \reck {#1}}{#2}} {#3}}{(\abse {#4}{#1}{#5})}}
\newcommand{\appe}[2]{\termc{#1\,#2}}
\newcommand{\letu}[2]{\termc{\letk\,\unitk = #1\,\ink\,#2}}
\newcommand{\lete}[4]{\termc{\letk\,\langle #1, #2\rangle = #3\,\ink\,{#4}}}
\newcommand{\inje}[2]{\termc {{#1}\,{#2}}}
\newcommand{\paire}[2]{\termc{\langle #1, #2\rangle}}
\newcommand{\newe}[1]{\tappe{\newk}{\typec{#1}}}
\newcommand{\forke}[1]{\forkk\,\termc{#1}}
\newcommand{\casee}[2]{\casek\;\termc #1\;\withk\;\termc{#2}}
\newcommand{\caseofe}[6]{\casee{#1}{\{{#2}_{#5}\;\overline{#3}_{#5}\rightarrow{#4}_{#5}\}_{{#5}\in{#6}}}}
\newcommand{\matche}[2]{\matchk\;\termc #1\;\withk\;\termc{#2}}
\newcommand{\matchwithe}[6]{\termc{\matche{#1}{\{{#2}_{#5}\;{#3}_{#5}\rightarrow{#4}_{#5}\}_{{#5}\in{#6}}}}}
\newcommand{\selecte}[1]{\selectk\,\termc{#1}} % Use C_i alone
\newcommand{\sendee}[2]{\sendk\;\termc{#1}\;\termc{#2}}
\newcommand{\receivee}[1]{\receivek\;\termc{#1}}
\newcommand{\closee}[1]{\terminatek\;\termc{#1}}
\newcommand{\waite}[1]{\waitk\;\termc{#1}}
\newcommand{\tappe}[2]{\termc{#1[}\typec{#2}\termc{]}}

\newcommand{\ebind}[2]{\termc{{#1}\blk\colon\typec{#2}}}
\newcommand{\eubind}[2]{\termc{{#1}\blk{\colon^{\!\!\!\star}}\typec{#2}}} % unrestricted

\newcommand\Normal[2][\pm]{\operator{nrm}^{\blk #1}\blk(\typec{#2}\blk)}

\newcommand\FVT[1]{\operator{fv} ({#1})} % in types
\newcommand\FVE[1]{\FVT{#1}} % in terms (exps and procs)
\newcommand\FVL[1]{\FVT{#1}} % in labels

% PROCESSES
\newcommand{\expp}[1]{\termc{\langle{#1}\rangle}}
\newcommand{\PAR}{\mid}
\newcommand{\parp}[2]{\termc{#1 \PAR #2}}
\newcommand{\newp}[3]{\termc{(\nu{#1}{#2})\,{#3}}}

% LABELLED TRANSITIONS
\newcommand{\lts}[1]{\xrightarrow{#1}}
\newcommand{\trans}[3]{\termc{#1}\lts{#2}\termc{#3}}
\newcommand{\transCtx}[3]{{#1}\lts{#2}{#3}}
\newcommand{\valueprocess}[1]{\termc{#1} \; \textsf{value}}
\newcommand{\forkl}[1]{\syntax{fork}\,{#1}}
\newcommand{\sendl}[2]{{#1}!{#2}}
\newcommand{\receivel}[2]{{#1}?{#2}}
\newcommand{\openl}[1]{{#1}!}
\newcommand{\closel}[1]{{#1?}}
\newcommand{\newl}[3]{(\nu{#1}{#2})\colon{#3}}
\newcommand{\scopel}[4]{(\nu{#1}{#2}){#3}!{#4}}
% \newcommand{\scopell}[3]{(\nu {#1}{#2}) {#3}!{#1}}
% \newcommand{\scoperl}[3]{(\nu {#1}{#2}) {#3}!{#2}}
\newcommand{\parl}[2]{{#1}\PAR{#2}}

% Translation of CFST to AlgST
\newcommand{\cfsttoalgst}[1]{\blk{\llbracket {#1} \rrbracket}}
\newcommand{\cfsttoprot}[1]{\blk{\llparenthesis {#1} \rrparenthesis}}
\newcommand{\algsttocfst}[1]{\blk{\Lbag {#1} \Rbag}}
\newcommand{\emptyseq}[0]{\varepsilon}
\newcommand{\iso}[1]{\gray{{#1}}}
\newcommand{\isodecl}[1]{\iso{#1}\gray{\colon}}
\newcommand{\isodef}[1]{\iso{#1}\gray{\,=\,}}

% JUDGMENTS
% \newcommand\isExpr[4][\Delta]{\twoCtx{#1}{#2} \vdash \termc{#3} : \typec{#4}} %deprecated
\newcommand\hasType[3][\Delta]{#1 \vdash {\typec{#2}} : {\kindc{#3}}} %deprecated
\newcommand\isConv[4][\Delta]{#1 \vdash \typec{#2} \equiv \typec{#3} : {\kindc
    {#4}}}
%\newcommand{\isEquiv}[2]{\Normal[+]{#1} =_\alpha \Normal[+]{#2}}
\newcommand{\isEquiv}[2]{\typec{#1}\equiv_A\typec{#2}} % Abbreviated version
\newcommand{\isEquivCtx}[2]{{#1}\equiv_\alpha{#2}}
% \newcommand{\isEquivCtx}[2]{{#1}\equiv_a{#2}} % Abbreviated version
\newcommand\isProc[2][\Gamma]{{#1} \vdash \termc{#2}}

\newcommand\isCtx[3][\Delta]{{#1} \vdash {#2}  \Leftarrow {#3}}
\newcommand\twoCtx[2]{#1 \mid #2}
\newcommand\typeSynth[3][\Delta]{#1 \vdash {\typec{#2}} \Rightarrow {#3}}
\newcommand\typeAgainst[3][\Delta]{#1 \vdash {\typec{#2}} \Leftarrow {#3}}

\newcommand\expSynth[5][\Delta]{\twoCtx{#1}{#2} \vdash \termc{#3} \Rightarrow
  \twoCtx{\typec{#4}}{#5}}
\newcommand\expAgainst[5][\Delta]{\twoCtx{#1}{#2} \vdash {\termc{#3}} \Leftarrow
  \twoCtx{\typec{#4}}{#5}}


% Operators
\newcommand{\tsubs}[2]{\blk[\typec{#1}\blk/\typec{#2}\blk]} % term substitution
\newcommand{\esubs}[2]{\blk[\termc{#1}\blk/\termc{#2}\blk]} % expression substitution

\newcommand{\declrel}[2]{\emph{#1}\hfill\fbox{{#2}}}
\newcommand{\Empty}{\cdot}

\newcommand\isSubKind[2]{\kindc{#1} \le \kindc{#2}}
\newcommand{\typeofop}{\operator{typeof}}
\newcommand\typeof[1]{\typeofop\operator{(}\termc{#1}\blk{)}}

\newcommand{\spec}{\operator{spec}}

%% Rules and axioms
\newcommand{\infrule}[3]{\inferrule[#1]{#2}{#3}}
\newcommand{\axiom}[2]{\infrule{#1}{}{#2}}
\newcommand{\infderiv}[3]{\inferrule*[lab=#1]{#2}{#3}}

% RULE NAMES
\newcommand{\rulename}[2]{\textsc{{#1}-{#2}}\xspace}

\newcommand{\rulenametype}[1]{\rulename{T}{#1}}  % T- for type formation
\newcommand{\rulenametconv}[1]{\rulename{C}{#1}}  % C- for type conversion
\newcommand{\rulenameexp}[1]{\rulename{E}{#1}}   % E- for expressions

% T- rule names
\newcommand{\rulenametypeEndW}{\rulenametype{End?}}
\newcommand{\rulenametypeEndT}{\rulenametype{End!}}
% \newcommand{\rulenametypeEnd}{\rulenametype{End}}
\newcommand{\rulenametypeDot}{\rulenametype{Dot}}
\newcommand{\rulenametypeDotIn}{\rulenametype{In}}
\newcommand{\rulenametypeDotOut}{\rulenametype{Out}}
\newcommand{\rulenametypeMsgPos}{\rulenametype{MsgPos}}
\newcommand{\rulenametypeMsgNeg}{\rulenametype{MsgNeg}}
\newcommand{\rulenametypeMsg}{\rulenametype{Msg}}
\newcommand{\rulenametypeMsgP}{\rulenametype{Msg2}}
%\newcommand{\rulenametypeDualG}{\rulenametype{DualG}}
\newcommand{\rulenametypeDualS}{\rulenametype{Dual}}
\newcommand{\rulenametypeDual}{\rulenametype{Dual}}
\newcommand{\rulenametypeInt}{\rulenametype{Int}}
\newcommand{\rulenametypeBool}{\rulenametype{Bool}}
\newcommand{\rulenametypeUnit}{\rulenametype{Unit}}
\newcommand{\rulenametypeUp}{\rulenametype{Up}}
\newcommand{\rulenametypePlus}{\rulenametype{Plus}}
\newcommand{\rulenametypeSemi}{\rulenametype{Semi}}
\newcommand{\rulenametypeMinus}{\rulenametype{Minus}}
\newcommand{\rulenametypeAssume}{\rulenametype{Assume}}
\newcommand{\rulenametypeArrow}{\rulenametype{Arrow}}
\newcommand{\rulenametypeGuarded}{\rulenametype{Guarded}}
\newcommand{\rulenametypeProtocol}{\rulenametype{Protocol}}
\newcommand{\rulenametypeVar}{\rulenametype{Var}}
\newcommand{\rulenametypePoly}{\rulenametype{Poly}}
\newcommand{\rulenametypeSub}{\rulenametype{Sub}}
\newcommand{\rulenametypePair}{\rulenametype{Pair}}

% C- rule names
\newcommand{\rulenametconvReflex}{\rulenametconv{Refl}}
\newcommand{\rulenametconvSymm}{\rulenametconv{Sym}}
\newcommand{\rulenametconvTrans}{\rulenametconv{Trans}}
\newcommand{\rulenametconvSub}{\rulenametconv{Sub}}
\newcommand{\rulenametconvDot}{\rulenametconv{Dot}}
\newcommand{\rulenametconvDoubleS}{\rulenametconv{DualInv}}
\newcommand{\rulenametconvDoubleG}{\rulenametconv{DoubleDual2}}
\newcommand{\rulenametconvDualEndW}{\rulenametconv{DualEnd?}}
\newcommand{\rulenametconvDualEndT}{\rulenametconv{DualEnd!}}
% \newcommand{\rulenametconvDualEnd}{\rulenametconv{DualEnd}}
\newcommand{\rulenametconvDualDot}{\rulenametconv{DualDot}}
\newcommand{\rulenametconvDualOutM}{\rulenametconv{DualOut}}
\newcommand{\rulenametconvDualOutP}{\rulenametconv{DualOut2}}
\newcommand{\rulenametconvDualInM}{\rulenametconv{DualIn}}
\newcommand{\rulenametconvDualInP}{\rulenametconv{DualIn2}}
\newcommand{\rulenametconvDualProto}{\rulenametconv{Protocol}}
\newcommand{\rulenametconvAll}{\rulenametconv{All}}

\newcommand{\rulenametconvNegOut}{\rulenametconv{NegOut}}
\newcommand{\rulenametconvNegIn}{\rulenametconv{NegIn}}
\newcommand{\rulenametconvPosOut}{\rulenametconv{PosOut}}
\newcommand{\rulenametconvPosIn}{\rulenametconv{PosIn}}
\newcommand{\rulenametconvPos}{\rulenametconv{Pos}}
\newcommand{\rulenametconvNegInv}{\rulenametconv{NegInv}}
\newcommand{\rulenametconvSemiIn}{\rulenametconv{SemiIn}}
\newcommand{\rulenametconvSemiOut}{\rulenametconv{SemiOut}}

% E- rule names
\newcommand{\rulenameexpConst}{\rulenameexp{Const}}
\newcommand{\rulenameexpUVar}{\rulenameexp{Var$\star$}}
\newcommand{\rulenameexpVar}{\rulenameexp{Var}}
% \newcommand{\rulenameexpUnVar}{\rulenameexp{UnVar}}
\newcommand{\rulenameexpAbs}{\rulenameexp{Abs}}
\newcommand{\rulenameexpAbsn}{\rulenameexp{Abs'}}
% \newcommand{\rulenameexpUnAbs}{\rulenameexp{UnAbs}}
\newcommand{\rulenameexpRec}{\rulenameexp{Rec}}
\newcommand{\rulenameexpApp}{\rulenameexp{App}}
\newcommand{\rulenameexpAppn}{\rulenameexp{App'}}
\newcommand{\rulenameexpTAbs}{\rulenameexp{TAbs}}
\newcommand{\rulenameexpTApp}{\rulenameexp{TApp}}
\newcommand{\rulenameexpAppAll}{\rulenameexp{AppAll}}
\newcommand{\rulenameexpNew}{\rulenameexp{New}}
\newcommand{\rulenameexpPair}{\rulenameexp{Pair}}
\newcommand{\rulenameexpLetUnit}{\rulenameexp{Let*}}
\newcommand{\rulenameexpLet}{\rulenameexp{Let}}
\newcommand{\rulenameexpCtor}{\rulenameexp{Ctor}}

\newcommand{\rulenameexpSelect}{\rulenameexp{Select}}
\newcommand{\rulenameexpCase}{\rulenameexp{Case}}
\newcommand{\rulenameexpMatch}{\rulenameexp{Match}}
% \newcommand{\rulenameexpSel}{\rulenameexp{Sel}}
\newcommand{\rulenameexpCheck}{\rulenameexp{Check}}

% Writing
\newcommand{\ie}{i.e.,\xspace} % comma
\newcommand{\eg}{e.g.\xspace}  % no comma, according to the Handbook for Scholars
\newcommand{\cf}{cf.\xspace}
\newcommand{\Cf}{Cf.\xspace}
\newcommand{\etal}{et al.\xspace} % al abbreviates aliae, no need for am \emph,
% but you may added if you feel like, according to the Handbook for Scholars
\renewcommand{\bibsection}{\section*{References}}

% listings

\lstset{
  captionpos=b,
  language=Haskell,
  basicstyle=\ttfamily\small,
  keywordstyle=\color{RedOrange}\ttfamily,
  commentstyle=\color{blue!50}\ttfamily,
  emphstyle=\color{RoyalBlue}\ttfamily, % types
  deletekeywords={repeat,List,Int,String,Char,Maybe,Either,Just,Nothing,Left,Right,seq,either,maybe,flip},
  morekeywords=[1]{protocol,select,send,receive,match,with,new,fork,wait,terminate},
  emph={Dual,Unit,Skip,Char,Int,String},
  numbers=none,
  keepspaces=true,
  literate=
  %{^}{$\uparrow$}1
  {^-1}{${}^{-1}$}2
  {|-}{$\vdash$}2
  {-|}{$\dashv$}2
  {->}{$\rightarrow$\,}2
  {|>}{$\rhd$}1
  {-o}{$\multimap$}2
  {=>}{$\Rightarrow$}2
  {[|}{$\llbracket$}1
  {|]}{$\rrbracket$}1
  {forall}{$\forall$}1
  {Lambda}{$\Lambda$}1
  {lambda}{$\lambda$}1
  {theta}{$\theta$}1
  {delta}{$\delta$}1
  {mu}{$\mu$}1
  {alpha}{$\alpha$}1
  {beta}{$\beta$}1
  {gamma}{$\gamma$}1
  {oplus}{$\oplus$}1
  {times}{$\times$}1
  {otimes}{$\otimes$}1
  {sigma}{$\sigma$}1
  {+\{}{$\oplus$\{}2
  {:P}{:\kindc{P}}2
  {:SL}{:\kindc{S}}2
  {:S}{:\kindc{S}}2
  {:M}{:\kindc{M}}2
  {:T}{:\kindc{T}}2
  {KP}{\kindc{P}}1
  {KS}{\kindc{S}}1
  {KM}{\kindc{M}}1
  {KT}{\kindc{T}}1
  {EndW}{\TEnd}3
  {EndT}{\TEnd}3
}

\mprset{lineskip=0ex}

% notes
\newcommand{\todo}[2][TODO]{[{\color{red}{#1}}: {#2}]}
\newcommand{\pt}[1]{\todo[pt]{#1}}
\newcommand{\js}[1]{\todo[js]{#1}}

% MW
\newcommand{\hl}[1]{\colorbox{violet!15}{$#1$}}
\newcommand{\hlt}[1]{\colorbox{violet!15}{#1}}
\newcommand{\subt}{\le}
%%% Local Variables:
%%% mode: latex
%%% TeX-master: "main"
%%% End:
